% Generated by analysis/apa_refs.py from refs.bib (APA 7). Do not edit by hand.
\begin{thebibliography}{12}
\bibitem[{Fastl \apaand{} Zwicker(2007)Fastl and Zwicker}]{zwicker2007}
Fastl, H., \& Zwicker, E. (2007). \emph{Psychoacoustics: Facts and models} (3rd ed.). Springer. \url{https://doi.org/10.1007/978-3-540-68888-4}

\bibitem[{Fletcher(1964)Fletcher}]{fletcher1964}
Fletcher, H. (1964). Normal vibration frequencies of a stiff piano string. \emph{The Journal of the Acoustical Society of America}, \emph{36}(1), 203--209. \url{https://doi.org/10.1121/1.1918933}

\bibitem[{Glasberg \apaand{} Moore(1990)Glasberg and Moore}]{glasberg1990}
Glasberg, B. R., \& Moore, B. C. J. (1990). Derivation of auditory filter shapes from notched-noise data. \emph{Hearing Research}, \emph{47}(1--2), 103--138. \url{https://doi.org/10.1016/0378-5955(90)90170-T}

\bibitem[{Marjieh et~al.(2024)Marjieh, Harrison, Lee, Deligiannaki, and Jacoby}]{marjieh2024}
Marjieh, R., Harrison, P. M. C., Lee, H., Deligiannaki, F., \& Jacoby, N. (2024). Timbral effects on consonance disentangle psychoacoustic mechanisms and suggest perceptual origins for musical scales. \emph{Nature Communications}, \emph{15}, 1482. \url{https://doi.org/10.1038/s41467-024-45812-z}

\bibitem[{Mashinter(2006)Mashinter}]{mashinter2006}
Mashinter, K. (2006). Calculating sensory dissonance: Some discrepancies arising from the models of Kameoka \& Kuriyagawa, and Hutchinson \& Knopoff. \emph{Empirical Musicology Review}, \emph{1}(2), 65--84. \url{https://doi.org/10.18061/1811/24077}

\bibitem[{Philharmonia Orchestra(n.d.)Philharmonia Orchestra}]{philharmonia}
Philharmonia Orchestra. (n.d.). \emph{Sound samples}. Retrieved August 2026, from \url{https://philharmonia.co.uk/resources/sound-samples/}

\bibitem[{Schlecht(2020)Schlecht}]{sethmatlab}
Schlecht, S. J. (2020). \emph{dissonanceSethares: MATLAB implementation of Sethares' dissonance measure} [Computer software]. \url{https://github.com/SebastianJiroSchlecht/dissonanceSethares}

\bibitem[{Sethares(1993)Sethares}]{sethares1993}
Sethares, W. A. (1993). Local consonance and the relationship between timbre and scale. \emph{The Journal of the Acoustical Society of America}, \emph{94}(3), 1218--1228. \url{https://doi.org/10.1121/1.408175}

\bibitem[{Sethares(2005)Sethares}]{sethares2005}
Sethares, W. A. (2005). \emph{Tuning, timbre, spectrum, scale} (2nd ed.). Springer. \url{https://doi.org/10.1007/b138848}

\bibitem[{Shamma \apaand{} Klein(2000)Shamma and Klein}]{shamma2000}
Shamma, S., \& Klein, D. (2000). The case of the missing pitch templates: How harmonic templates emerge in the early auditory system. \emph{The Journal of the Acoustical Society of America}, \emph{107}(5), 2631--2644. \url{https://doi.org/10.1121/1.428649}

\bibitem[{University of Iowa Electronic Music Studios(n.d.)University of Iowa Electronic Music Studios}]{iowamis}
University of Iowa Electronic Music Studios. (n.d.). \emph{Musical instrument samples}. Retrieved August 2026, from \url{https://theremin.music.uiowa.edu/MIS.html}

\bibitem[{Vassilakis(2001)Vassilakis}]{vassilakis2001}
Vassilakis, P. N. (2001). \emph{Perceptual and physical properties of amplitude fluctuation and their musical significance} [Doctoral dissertation, University of California, Los Angeles].

\end{thebibliography}
